\chapter{Methods}
\label{cha:methods}

\section{Core Idea}

\newcommand{\R}{\mathbb{R}}
\newcommand{\E}{\mathbb{E}}
\newcommand{\PG}{\mathrm{PG}}
\newcommand{\xo}{x^o}
\newcommand{\xm}{x^m}

\subsection{Setup} Consider a logistic regression, with missing features in the observation.\\ 
Let $x := (\xo, \xm) \in \R^{d}$. Then the logistic regression is defined as
\[
  p(y \mid x) := \sigma (w^T x), \qquad w \in \R^{d}.
\]
Additionally we have a PC that models the density $p(x)$.

\subsection{Goal} Compute the prediction in a probabilistically principled manner using
\[
  p(y \mid \xo) := \E_{p(\xm \mid \xo)} \left[\sigma \left( w^T x \right) \right].
\]

If it were not for the non-linear sigmoid function $\sigma$, we would be able to calculate this in a smooth 
and decomposable PC, by pulling the Expectation -which is linear- into the leaf nodes. To tackle this we propose to use two tricks.

\subsection{Trick 1: Pólya-Gamma}

To replace the sigmoid function we use the PG-Sigmoid identity, shown in \cite{polyagamma}, that states:

\[
  \sigma(z) = \frac{1}{2} \exp \left(\frac{1}{2}z\right) \
  \E_{\gamma \sim \PG(1,0)} \left[\exp \left(-\frac{1}{2} \gamma z^2 \right)\right]
\]
for $z \in \R$.

If we set $z := w^T x$, and insert into our expectation, we get:
\begin{align*}
  p(y \mid \xo)
  & = \E_{p(\xm \mid \xo)} \left[\sigma(z)\right] \\
  & = \E_{p(\xm \mid \xo)} \left[
       \frac{1}{2}\exp \left(\frac{1}{2}z\right)
       \E_{\gamma \sim \PG(1,0)} \left[\exp \left(-\frac{1}{2}\gamma z^2 \right)\right]
     \right] \\
  & = \frac{1}{2} \E_{p(\xm \mid \xo)} \E_{\gamma \sim \PG(1,0)}
    \left[ \exp \left(\frac{1}{2}z - \frac{1}{2}\gamma z^2 \right) \right]
\end{align*}

Now, we have a term without a sigmoid, but the $x$ still appears in a quadratic form.

\subsection{Trick 2: Hubbard-Stratonovich (HS)}

To get rid of the quadratic form, we apply the Hubbard--Stratonovich identity \cite{stratonovich1957, hubbard1959}, that states:

\begin{equation}
    \exp \left( -\tfrac{1}{2} a h^2 \right) = \sqrt{  \frac{1}{2 \pi a}} \int_{-\infty}^{\infty} \exp \left(-\frac{u^2}{2a} -i h u\right) \mathrm{d}u,
\end{equation}

In plain words, it makes a quadratic term linear at the cost of introducing a helper variable that we need to integrate over.

Currently, in our term, $x$ appears quadratically as well as linearly. To be able to apply HS, we complete the square (again using $z = w^T x$):

\begin{align}
  \frac{1}{2}z - \frac{1}{2}\gamma z^2
  &= -\frac{1}{2}\gamma \left(z^2 - \frac{z}{\gamma}\right) \notag\\
  &= -\frac{1}{2}\gamma \left(\left(z - \frac{1}{2\gamma}\right)^{2} - \frac{1}{4\gamma^2}\right) \notag\\
  &= -\frac{1}{2}\gamma \left(z - \frac{1}{2\gamma}\right)^{2} + \frac{1}{2}\gamma \cdot \frac{1}{4\gamma^2} \notag\\
  &= -\frac{1}{2}\gamma \left(z - \frac{1}{2\gamma}\right)^{2} + \frac{1}{8\gamma}.
\end{align}

Now we apply HS with $a = \gamma$ and $h = z - \frac{1}{2\gamma}$:

\begin{align}
  \exp \left(\frac{1}{2}z - \frac{1}{2}\gamma z^2 \right)
  &= \exp \left(\frac{1}{8\gamma}\right) \exp \left(-\frac{1}{2} \gamma \left(z - \frac{1}{2\gamma}\right)^{2} \right) \notag \\
  &= \exp \left(\frac{1}{8\gamma}\right) \sqrt{\frac{1}{2\pi\gamma}} \int_{-\infty}^{\infty} \exp \left(-\frac{u^2}{2\gamma} - i \left(z - \frac{1}{2\gamma}\right) u\right) \mathrm{d}u\notag \\
  &= \exp \left(\frac{1}{8\gamma}\right) \int_{-\infty}^{\infty} \underbrace{\frac{1}{\sqrt{2\pi\gamma}} \exp \left(-\frac{u^2}{2\gamma}\right)}_{\mathcal{N}(u; 0, \gamma)} \exp \left(-i \left(z - \frac{1}{2\gamma}\right) u \right) \mathrm{d}u \notag \\
  &= \exp \left(\frac{1}{8\gamma}\right) \E_{u \sim \mathcal{N}(0, \gamma)} \left[ \exp \left(-i \left(z - \frac{1}{2\gamma}\right) u \right) \right] \notag \\
  &= \exp \left(\frac{1}{8\gamma}\right) \E_{u \sim \mathcal{N}(0, \gamma)} \left[ \exp \left(\frac{iu}{2\gamma} - izu \right) \right]
\end{align}

Substituting back into $p(y \mid \xo)$, we get:

\begin{align*}
  p(y \mid \xo)
  &= \frac{1}{2} \E_{p(\xm \mid \xo)} \E_{\gamma \sim \PG(1,0)}
     \left[ \exp \left(\frac{1}{2}z - \frac{1}{2}\gamma z^2 \right) \right] \\
  &= \frac{1}{2} \E_{p(\xm \mid \xo)} \E_{\gamma \sim \PG(1,0)}
     \left[ \exp \left(\frac{1}{8\gamma}\right) \E_{u \sim \mathcal{N}(0, \gamma)}
     \left[ \exp \left(\frac{iu}{2\gamma} - izu \right) \right] \right] \\
  &= \frac{1}{2} \E_{\gamma \sim \PG(1,0)} \E_{u \sim \mathcal{N}(0, \gamma)}
     \left[ \exp \left(\frac{1 + 4iu}{8\gamma}\right)
     \E_{p(\xm \mid \xo)} \left[ \exp \left(-izu \right) \right] \right]
\end{align*}

\subsection{Characteristic Function}

The characteristic function of a random variable $X$ is defined as $\phi_X(s) = \E \left[ \exp(isX) \right]$.
The innermost expectation is only over $\xm$, with $\xo$ fixed, so we split $z = w^T x$ into its observed and missing parts:
\[
  z = w^T x = w_o^T \xo + w_m^T \xm.
\]
The $\xo$ part is constant under $\E_{p(\xm \mid \xo)}$ and factors out:
\begin{align*}
  \E_{p(\xm \mid \xo)} \left[ \exp(-izu) \right]
  &= \E_{p(\xm \mid \xo)} \left[ \exp\left(-i(w_o^T \xo + w_m^T \xm)u\right) \right] \\
  &= \exp\left(-iw_o^T \xo u\right) \cdot \E_{p(\xm \mid \xo)} \left[ \exp\left(-iw_m^T \xm u\right) \right] \\
  &= \exp\left(-iw_o^T \xo u\right) \cdot \phi_{\xm \mid \xo}(-uw_m)
\end{align*}
So the input to the characteristic function is $s := -uw_m$.

\subsection{Full Formula}

Plugging the result of the previous subsection back into the outer
expression for $p(y \mid \xo)$, we obtain:
\begin{align*}
  p(y \mid \xo)
  &= \frac{1}{2} \E_{\gamma \sim \PG(1,0)} \E_{u \sim \mathcal{N}(0, \gamma)}
     \left[ \exp \left(\frac{1 + 4iu}{8\gamma}\right)
     \cdot \exp\left(-iw_o^T \xo u\right) \cdot \phi_{\xm \mid \xo}(-uw_m) \right] \\
  &= \frac{1}{2} \E_{\gamma \sim \PG(1,0)} \E_{u \sim \mathcal{N}(0, \gamma)}
     \left[ \exp \left(\frac{1}{8\gamma} + \frac{iu}{2\gamma} - iw_o^T \xo u\right)
     \cdot \phi_{\xm \mid \xo}(-uw_m) \right] \\
  &= \frac{1}{2} \E_{\gamma \sim \PG(1,0)} \E_{u \sim \mathcal{N}(0, \gamma)}
     \left[ \exp \left(\frac{1}{8\gamma} + iu \left(\frac{1}{2\gamma} - w_o^T \xo\right)\right)
     \cdot \phi_{\xm \mid \xo}(-uw_m) \right]
\end{align*}

We propose this as a final expression that we can tractably evaluate.

\section{Implementation}

Two ingredients are still abstract in the formula above: the conditional
characteristic function $\phi_{\xm \mid \xo}$, which depends on the
density model, and the inner $\mathcal{N}(0, \gamma)$ expectation, which
needs to be approximated. This section makes both concrete.

\subsection{Conditional CF for a GMM}

For a Gaussian mixture model $p(x) = \sum_{k=1}^{K} \pi_k \,
\mathcal{N}(x; \mu_k, \Sigma_k)$ with diagonal covariances $\Sigma_k =
\mathrm{diag}(\sigma_k^2)$, the conditional density $p(\xm \mid \xo)$ is
again a mixture, with the same component shapes restricted to the missing
coordinates and reweighted by the posterior responsibility of each
component given $\xo$:
\[
  \tilde{\pi}_k(\xo) = \frac{\pi_k \, \mathcal{N}(\xo; \mu_k^o, \Sigma_k^o)}{\sum_{j=1}^{K} \pi_j \, \mathcal{N}(\xo; \mu_j^o, \Sigma_j^o)}.
\]
Since the characteristic function of $\mathcal{N}(\mu, \Sigma)$ at $s$ is
$\exp \left( i s^T \mu - \tfrac{1}{2} s^T \Sigma s \right)$, and the CF of
a mixture is the same convex combination of component CFs, we get a
closed-form expression:
\[
  \phi_{\xm \mid \xo}(s) = \sum_{k=1}^{K} \tilde{\pi}_k(\xo) \,
  \exp \left( i s^T \mu_k^m - \tfrac{1}{2} s^T \Sigma_k^m s \right),
\]
where $\mu_k^m, \Sigma_k^m$ are the component means and covariances
restricted to the missing coordinates.

\subsection{Conditional CF for a General PC}

More generally, any smooth and decomposable probabilistic circuit admits
the conditional CF in closed form by recursive evaluation, since
$\phi_X(s) = \E[\exp(is^T X)]$ is itself an expectation of a function
that factorises over disjoint variable scopes, and PCs are by
construction tractable for exactly this kind of expectation. Concretely,
for a node $n$ over scope $\mathrm{sc}(n)$ with input $s$ restricted to
that scope:
\begin{itemize}
  \item \textbf{Leaf node} over a single variable $X$ with density
        $p_n(x)$: $\phi_n(s) = \int \exp(isx) \, p_n(x) \, \mathrm{d}x$, which
        is closed-form for the standard leaf families (Gaussian, categorical,
        \dots).
  \item \textbf{Product node} with children $n_1, \dots, n_C$ on disjoint
        scopes (decomposability): variables in different children are
        independent, so the CF factorises,
        $\phi_n(s) = \prod_{c=1}^{C} \phi_{n_c}(s_{\mathrm{sc}(n_c)})$.
  \item \textbf{Sum node} with weights $w_c \geq 0$, $\sum_c w_c = 1$:
        the CF is the same convex combination of children's CFs,
        $\phi_n(s) = \sum_{c=1}^{C} w_c \, \phi_{n_c}(s)$.
\end{itemize}
Conditioning on $\xo$ amounts to evaluating the observed-scope leaves at
$\xo$ and renormalising the sum weights to posterior responsibilities, in
exact analogy to the GMM case above; the missing-scope evaluation then
returns $\phi_{\xm \mid \xo}(s)$ in a single bottom-up pass.

\subsection{Gauss--Hermite Quadrature}
The inner Gaussian expectation is amenable to Gauss--Hermite quadrature, with
deterministic nodes $\{t_i\}_{i=1}^{N}$ and weights $\{\omega_i\}_{i=1}^{N}$
such that

\[
  \int_{-\infty}^{\infty} e^{-t^2} g(t) \, \mathrm{d}t \approx \sum_{i=1}^{N} \omega_i \, g(t_i)
\]

To use this we need to bring the $\mathcal{N}(0, \gamma)$
expectation into the form $\int e^{-t^2} g(t)\,\mathrm{d}t$. Writing the
expectation out as an integral against the Gaussian density gives

\[
  \E_{u \sim \mathcal{N}(0, \gamma)}[f(u)]
  = \int_{-\infty}^{\infty} \frac{1}{\sqrt{2\pi\gamma}} \exp \!\left(-\frac{u^2}{2\gamma}\right) f(u) \, \mathrm{d}u.
\]

Apply the substitution $u = \sqrt{2\gamma}\,t$. It changes three pieces of
the integrand:
\begin{itemize}
  \item the differential: $\mathrm{d}u = \sqrt{2\gamma}\,\mathrm{d}t$,
  \item the Gaussian exponent: $\dfrac{u^2}{2\gamma} = \dfrac{(\sqrt{2\gamma}\,t)^2}{2\gamma} = \dfrac{2\gamma\, t^2}{2\gamma} = t^2$,
  \item the argument of $f$: $f(u) = f \!\left(\sqrt{2\gamma}\,t\right)$,
\end{itemize}
and the integration limits are unchanged since $\sqrt{2\gamma} > 0$. Plugging
all three substitutions into the integral:
\begin{align*}
  \E_{u \sim \mathcal{N}(0, \gamma)}[f(u)]
  &= \int_{-\infty}^{\infty} \frac{1}{\sqrt{2\pi\gamma}} \, e^{-t^2} f \!\left(\sqrt{2\gamma}\,t\right) \sqrt{2\gamma} \, \mathrm{d}t \\
  &= \frac{\sqrt{2\gamma}}{\sqrt{2\pi\gamma}} \int_{-\infty}^{\infty} e^{-t^2} f \!\left(\sqrt{2\gamma}\,t\right) \mathrm{d}t.
\end{align*}
The constant prefactor simplifies to
\[
  \frac{\sqrt{2\gamma}}{\sqrt{2\pi\gamma}}
  = \sqrt{\frac{2\gamma}{2\pi\gamma}}
  = \sqrt{\frac{1}{\pi}}
  = \frac{1}{\sqrt{\pi}},
\]
so the Gaussian expectation becomes
\[
  \E_{u \sim \mathcal{N}(0, \gamma)}[f(u)]
  = \frac{1}{\sqrt{\pi}} \int_{-\infty}^{\infty} e^{-t^2} f \!\left(\sqrt{2\gamma}\,t\right) \mathrm{d}t.
\]

Which we then can apply by
\[
  \E_{u \sim \mathcal{N}(0, \gamma)}[f(u)]
  \approx \frac{1}{\sqrt{\pi}} \sum_{i=1}^{N} \omega_i\, f \!\left(\sqrt{2\gamma}\,t_i\right).
\]

Plugging this into the inner expectation of our final formula, with\\
\[
  f(u) = \exp \!\left( \frac{1}{8\gamma} + i u \left(\frac{1}{2\gamma} - w_o^T \xo \right) \right) \phi_{\xm \mid \xo}(-u w_m)
\]
gives

\[
  p(y \mid \xo) \approx
  \frac{1}{2\sqrt{\pi}} \E_{\gamma \sim \PG(1,0)} \left[
    \sum_{i=1}^{N} \omega_i \,
    \exp \!\left( \frac{1}{8\gamma} + i \sqrt{2\gamma}\, t_i \left(\frac{1}{2\gamma} - w_o^T \xo \right) \right)
    \phi_{\xm \mid \xo}\!\left(-\sqrt{2\gamma}\, t_i\, w_m\right)
  \right]
\]

\subsection{Gauss--Pólya-Gamma Quadrature}
\label{subsec:pg-quadrature}

The outer expectation over $\gamma \sim \PG(1,0)$ can likewise be handled
deterministically by constructing a bespoke quadrature rule for the PG(1,0)
measure.  The key ingredient is the \emph{Golub--Welsch algorithm}
\cite{golubwelsch1969}, which produces $n$-point quadrature nodes and weights
from the first $2n$ moments of any positive measure.

\paragraph{Moments of PG$(1,0)$.}
The Laplace transform of $\PG(1,0)$ is known in closed form:
\[
  \mathcal{L}(t)
  := \E_{\gamma \sim \PG(1,0)}\!\left[e^{-t\gamma}\right]
  = \operatorname{sech}\!\left(\sqrt{\tfrac{t}{2}}\right)
  = \frac{1}{\cosh\!\left(\sqrt{t/2}\right)}.
\]
The raw moments follow from derivatives at zero:
\[
  \mu_k := \E[\gamma^k] = (-1)^k\, \mathcal{L}^{(k)}(0),
\]
which can be evaluated numerically to arbitrary precision
(we use \texttt{mpmath} at 200 decimal places to avoid catastrophic
cancellation for large $k$).

\paragraph{Golub--Welsch construction.}
Given $\mu_0, \dots, \mu_{2n-1}$, the algorithm proceeds as follows.
Form the $n \times n$ Hankel moment matrices
\[
  H_{ij} = \mu_{i+j}, \qquad H'_{ij} = \mu_{i+j+1}, \quad i,j = 0,\dots,n-1.
\]
Compute the Cholesky factor $H = R^\top R$ and the Jacobi (tridiagonal) matrix
\[
  J = R^{-\top} H' R^{-1}, \quad \text{symmetrised as } J \leftarrow \tfrac{1}{2}(J + J^\top).
\]
Eigendecompose $J = V \Lambda V^\top$.  Then the quadrature nodes and weights are
\[
  \gamma_j = \lambda_j, \qquad w_j = \mu_0\, v_{0j}^2, \quad j = 1,\dots,n,
\]
where $\lambda_j$ are the eigenvalues and $v_{0j}$ is the first component of
the $j$-th eigenvector.  The resulting rule integrates polynomials in $\gamma$
exactly up to degree $2n-1$.

\paragraph{Numerical considerations.}
The Hankel matrix $H$ becomes numerically singular in double precision for
$n \gtrsim 12$.  We therefore perform the entire Golub--Welsch computation in
high-precision arithmetic (\texttt{mpmath}, \texttt{dps}$=200$) and convert
to float\,32 only at the end.  The maximum reliable order is $n = 19$; beyond
that, eigenvalues become negative at any tested precision.

\paragraph{Fully deterministic formula.}
Combining Gauss--PG nodes $(\gamma_j, w_j)_{j=1}^{M}$ with Gauss--Hermite
nodes $(t_i, \omega_i)_{i=1}^{N}$ and applying both quadrature rules to the
full double expectation yields the \emph{zero-variance} approximation
\[
  p(y \mid \xo) \approx
  \frac{1}{2\sqrt{\pi}} \sum_{j=1}^{M} w_j \sum_{i=1}^{N} \omega_i \,
  \exp\!\left(\frac{1}{8\gamma_j} + i\sqrt{2\gamma_j}\,t_i
    \left(\frac{1}{2\gamma_j} - w_o^T \xo\right)\right)
  \phi_{\xm \mid \xo}\!\left(-\sqrt{2\gamma_j}\,t_i\,w_m\right).
\]

\subsection{Algorithms}

Putting the pieces together for a diagonal-covariance GMM density.
Algorithm~\ref{alg:pcpg-mc} uses Monte Carlo for both integrals.
Algorithm~\ref{alg:pcpg-gh} keeps MC for the outer PG integral but replaces
the inner Gaussian MC with Gauss--Hermite quadrature, eliminating inner
variance.
Algorithm~\ref{alg:pcpg-quad} further replaces the PG Monte Carlo with
Gauss--PG quadrature, making the estimator fully deterministic.

\begin{algorithm}[H]
\caption{PCPG: expected sigmoid under a GMM, full Monte Carlo}
\label{alg:pcpg-mc}
\begin{algorithmic}[1]
\Require LR weights $w$, partial observation $\xo$ (missing entries marked),
GMM parameters $\{\pi_k, \mu_k, \Sigma_k\}_{k=1}^{K}$,
PG samples $M$, Gaussian samples $L$
\State Split $w = (w_o, w_m)$ and $\mu_k = (\mu_k^o, \mu_k^m)$, $\Sigma_k = (\Sigma_k^o, \Sigma_k^m)$ by the observed/missing mask
\State $\tilde{\pi}_k \gets \dfrac{\pi_k \, \mathcal{N}(\xo;\mu_k^o,\Sigma_k^o)}{\sum_{j}\pi_j \,\mathcal{N}(\xo;\mu_j^o,\Sigma_j^o)}$ \hfill for $k=1,\dots,K$
\State $c \gets w_o^T \xo$
\State $S \gets 0$
\For{$j = 1, \dots, M$}
  \State Sample $\gamma_j \sim \PG(1, 0)$
  \For{$\ell = 1, \dots, L$}
    \State Sample $u \sim \mathcal{N}(0, \gamma_j)$
    \State $s \gets -u\, w_m$
    \State $\phi \gets \sum_{k=1}^{K} \tilde{\pi}_k \exp \!\left( i s^T \mu_k^m - \tfrac{1}{2} s^T \Sigma_k^m s \right)$
    \State $\mathrm{factor} \gets \exp \!\left( \dfrac{1}{8\gamma_j} + i u \left( \dfrac{1}{2\gamma_j} - c \right) \right)$
    \State $S \gets S + \mathrm{factor} \cdot \phi$
  \EndFor
\EndFor
\State \Return $\dfrac{1}{2\,M\,L} \cdot \mathrm{Re}(S)$
\end{algorithmic}
\end{algorithm}

\begin{algorithm}[H]
\caption{PCPG-GH: expected sigmoid under a GMM, MC outer + Gauss--Hermite inner}
\label{alg:pcpg-gh}
\begin{algorithmic}[1]
\Require LR weights $w$, partial observation $\xo$ (missing entries marked),
GMM parameters $\{\pi_k, \mu_k, \Sigma_k\}_{k=1}^{K}$,
PG samples $M$, Hermite order $N$
\State $(t_i, \omega_i)_{i=1}^{N} \gets \texttt{hermgauss}(N)$
\Comment{deterministic nodes/weights, precomputed once}
\State Split $w = (w_o, w_m)$ and $\mu_k = (\mu_k^o, \mu_k^m)$, $\Sigma_k = (\Sigma_k^o, \Sigma_k^m)$ by the observed/missing mask
\State $\tilde{\pi}_k \gets \dfrac{\pi_k \, \mathcal{N}(\xo;\mu_k^o,\Sigma_k^o)}{\sum_{j}\pi_j \,\mathcal{N}(\xo;\mu_j^o,\Sigma_j^o)}$ \hfill for $k=1,\dots,K$
\State $c \gets w_o^T \xo$
\State $S \gets 0$
\For{$j = 1, \dots, M$}
  \State Sample $\gamma_j \sim \PG(1, 0)$
  \For{$i = 1, \dots, N$}
    \State $u \gets \sqrt{2\gamma_j}\, t_i$
    \State $s \gets -u\, w_m$
    \State $\phi \gets \sum_{k=1}^{K} \tilde{\pi}_k \exp \!\left( i s^T \mu_k^m - \tfrac{1}{2} s^T \Sigma_k^m s \right)$
    \State $\mathrm{factor} \gets \exp \!\left( \dfrac{1}{8\gamma_j} + i u \left( \dfrac{1}{2\gamma_j} - c \right) \right)$
    \State $S \gets S + \omega_i \cdot \mathrm{factor} \cdot \phi$
  \EndFor
\EndFor
\State \Return $\dfrac{1}{2\sqrt{\pi}\,M} \cdot \mathrm{Re}(S)$
\end{algorithmic}
\end{algorithm}

\begin{algorithm}[H]
\caption{PCPG-Quad: fully deterministic expected sigmoid under a GMM}
\label{alg:pcpg-quad}
\begin{algorithmic}[1]
\Require LR weights $w$, partial observation $\xo$ (missing entries marked),
GMM parameters $\{\pi_k, \mu_k, \Sigma_k\}_{k=1}^{K}$,
PG quadrature order $M$ ($M \leq 19$), Hermite order $N$
\State $(\gamma_j, w_j)_{j=1}^{M} \gets \texttt{pg\_gauss\_quadrature}(M)$
\Comment{Golub--Welsch, computed once at high precision; see \S\ref{subsec:pg-quadrature}}
\State $(t_i, \omega_i)_{i=1}^{N} \gets \texttt{hermgauss}(N)$
\State Split $w = (w_o, w_m)$ and $\mu_k = (\mu_k^o, \mu_k^m)$, $\Sigma_k = (\Sigma_k^o, \Sigma_k^m)$ by the observed/missing mask
\State $\tilde{\pi}_k \gets \dfrac{\pi_k \, \mathcal{N}(\xo;\mu_k^o,\Sigma_k^o)}{\sum_{j}\pi_j \,\mathcal{N}(\xo;\mu_j^o,\Sigma_j^o)}$ \hfill for $k=1,\dots,K$
\State $c \gets w_o^T \xo$
\State $S \gets 0$
\For{$j = 1, \dots, M$}
  \For{$i = 1, \dots, N$}
    \State $u \gets \sqrt{2\gamma_j}\, t_i$
    \State $s \gets -u\, w_m$
    \State $\phi \gets \sum_{k=1}^{K} \tilde{\pi}_k \exp \!\left( i s^T \mu_k^m - \tfrac{1}{2} s^T \Sigma_k^m s \right)$
    \State $\mathrm{factor} \gets \exp \!\left( \dfrac{1}{8\gamma_j} + i u \left( \dfrac{1}{2\gamma_j} - c \right) \right)$
    \State $S \gets S + w_j \cdot \omega_i \cdot \mathrm{factor} \cdot \phi$
  \EndFor
\EndFor
\State \Return $\dfrac{1}{2\sqrt{\pi}} \cdot \mathrm{Re}(S)$
\end{algorithmic}
\end{algorithm}

\section{Discussion}

\subsection{Variance of the MC Estimator}

The Monte Carlo estimator draws $N$ samples $x^{(i)} \sim p(\xm \mid \xo)$ and returns their average sigmoid value:
\[
  \hat{I}_\text{MC} = \frac{1}{N} \sum_{i=1}^{N} \sigma\!\left(w^T x^{(i)}\right).
\]
This is a standard sample mean, so its variance is
\[
  \operatorname{Var}(\hat{I}_\text{MC}) = \frac{\sigma^2_\text{MC}}{N}, \qquad
  \sigma^2_\text{MC} = \operatorname{Var}_{x \sim p(x \mid \xo)}\!\left[\sigma(w^T x)\right].
\]
Since $\sigma(w^T x) \in (0, 1)$, the variance of any $[0,1]$-valued random variable is at most $\tfrac{1}{4}$ (attained by a Bernoulli$(\tfrac{1}{2})$ distribution), giving the tight upper bound
\[
  \operatorname{Var}(\hat{I}_\text{MC}) \leq \frac{1}{4N}.
\]

\subsection{Variance of the PCPG Estimator}

PCPG is a two-level Monte Carlo estimator. Let $h(\gamma, u)$ denote the integrand
\[
  h(\gamma, u) := \operatorname{Re}\!\left[
    \exp\!\left(\frac{1}{8\gamma} + iu\!\left(\frac{1}{2\gamma} - w_o^T \xo\right)\right)
    \phi_{\xm \mid \xo}(-u\,w_m)
  \right],
\]
and let $g(\gamma) := \E_{u \sim \mathcal{N}(0,\gamma)}[h(\gamma, u)]$ be the inner expectation.
The estimator returns
\[
  \hat{I}_\text{PCPG} = \frac{1}{2} \cdot \frac{1}{n_{pg}} \sum_{j=1}^{n_{pg}}
  \underbrace{\frac{1}{n_\text{gauss}} \sum_{k=1}^{n_\text{gauss}} h\!\left(\gamma^{(j)}, u^{(k,j)}\right)}_{\approx\; g(\gamma^{(j)})}.
\]
Applying the law of total variance to the inner mean gives
\[
  \operatorname{Var}(\hat{I}_\text{PCPG})
  = \frac{1}{4}\left[
      \frac{\sigma^2_B}{n_{pg}} + \frac{\sigma^2_A}{n_{pg} \cdot n_\text{gauss}}
    \right],
\]
where
\begin{align*}
  \sigma^2_B &= \operatorname{Var}_{\gamma \sim \PG(1,0)}\!\left[g(\gamma)\right],
  \quad \text{(outer PG variance, irreducible by increasing } n_\text{gauss}\text{)}\\
  \sigma^2_A &= \E_{\gamma \sim \PG(1,0)}\!\left[\operatorname{Var}_{u \sim \mathcal{N}(0,\gamma)}\!\left[h(\gamma,u)\right]\right].
  \quad \text{(inner Gaussian variance)}
\end{align*}
Unlike the MC case, neither $\sigma^2_A$ nor $\sigma^2_B$ is bounded by $\tfrac{1}{4}$.
The integrand $h(\gamma, u)$ contains the factor $\exp(1/(8\gamma))$, which grows
without bound as $\gamma \to 0^+$; since PG$(1,0)$ assigns positive probability to
arbitrarily small $\gamma$, both variance components are in principle unbounded.

\subsection{Variance of the PCPG-GH Estimator}

PCPG-GH (Algorithm~\ref{alg:pcpg-gh}) replaces the inner Monte Carlo sum with
Gauss--Hermite quadrature.  Using $N_H$ Hermite nodes $(t_i, \omega_i)$, the inner
integral is approximated deterministically as
\[
  \tilde{g}(\gamma) := \frac{1}{\sqrt{\pi}} \sum_{i=1}^{N_H} \omega_i\,
  h\!\left(\gamma,\, \sqrt{2\gamma}\,t_i\right) \;\approx\; g(\gamma).
\]
The estimator is then a \emph{single-level} Monte Carlo average over the PG draws:
\[
  \hat{I}_\text{PCPG-GH} = \frac{1}{2} \cdot \frac{1}{n_{pg}} \sum_{j=1}^{n_{pg}} \tilde{g}\!\left(\gamma^{(j)}\right).
\]
There is no inner randomness, so the law of total variance reduces immediately to
\[
  \operatorname{Var}(\hat{I}_\text{PCPG-GH})
  = \frac{1}{4} \cdot \frac{\operatorname{Var}_\gamma[\tilde{g}(\gamma)]}{n_{pg}}
  \;\xrightarrow{N_H \to \infty}\;
  \frac{\sigma^2_B}{4\,n_{pg}}.
\]
The $\sigma^2_A$ term present in PCPG is entirely eliminated: Gauss--Hermite
quadrature removes all inner variance at the cost of $N_H$ deterministic CF
evaluations per PG draw instead of $n_\text{gauss}$ random ones.

\paragraph{Bias of PCPG-GH.}
Unlike the full-MC PCPG estimator, PCPG-GH is \emph{biased} for any finite $N_H$
because $\tilde{g}(\gamma) \neq g(\gamma)$ in general:
\[
  \operatorname{Bias}(\hat{I}_\text{PCPG-GH})
  = \frac{1}{2}\,\E_{\gamma \sim \PG(1,0)}\!\left[\tilde{g}(\gamma) - g(\gamma)\right].
\]
The inner error $\tilde{g}(\gamma) - g(\gamma)$ is the error of an $N_H$-point
Gauss--Hermite rule applied to the integral $\int e^{-t^2} h(\gamma, \sqrt{2\gamma}\,t)\,\mathrm{d}t$.
The classical remainder formula states that for a $2N_H$-times differentiable
integrand $f$:
\[
  \int_{-\infty}^{\infty} e^{-t^2} f(t)\,\mathrm{d}t
  - \sum_{i=1}^{N_H} \omega_i f(t_i)
  = \frac{\sqrt{\pi}\,(N_H)!}{2^{2N_H}\,(2N_H)!}\,f^{(2N_H)}(\xi)
\]
for some $\xi \in \mathbb{R}$.
The prefactor $\frac{(N_H)!}{2^{2N_H}(2N_H)!}$ decays super-exponentially in
$N_H$ (faster than any polynomial), so the inner bias is negligible in practice
for $N_H \geq 20$: the integrand $h(\gamma, \sqrt{2\gamma}\,t)$ is analytic in
$t$ (it involves complex exponentials and the GMM CF, which are entire functions),
so all derivatives exist and are bounded on compact sets.
In summary, PCPG-GH trades a small, rapidly vanishing systematic error for the
complete elimination of inner-loop variance.

\subsection{Variance and Bias of the PCPG-Quad Estimator}

PCPG-Quad (Algorithm~\ref{alg:pcpg-quad}) is fully deterministic: both the outer
PG expectation and the inner Gaussian expectation are replaced by quadrature rules,
so it has zero variance but nonzero bias from both approximations.

\paragraph{Variance.}
Since the estimator is deterministic, $\operatorname{Var}(\hat{I}_\text{PCPG-Quad}) = 0$.
Running it multiple times with the same inputs always returns the same value.

\paragraph{Bias from the inner (Gauss--Hermite) quadrature.}
The inner approximation $\tilde{g}(\gamma) \approx g(\gamma)$ carries the same
super-exponentially vanishing error as in PCPG-GH, and is negligible for
$N_H \geq 20$.

\paragraph{Bias from the outer (Gauss--PG) quadrature.}
The outer integral $\frac{1}{2}\E_{\gamma \sim \PG(1,0)}[\tilde{g}(\gamma)]$ is
replaced by the $M$-point quadrature sum $\frac{1}{2}\sum_{j=1}^{M} w_j\,\tilde{g}(\gamma_j)$.
An $M$-point Gauss quadrature rule for measure $p_\text{PG}$ is exact for
polynomials in $\gamma$ up to degree $2M-1$.  The integrand $\tilde{g}(\gamma)$
is not a polynomial: it contains the factor $e^{1/(8\gamma)}$, which has an
essential singularity at $\gamma = 0$.  Consequently the quadrature error does
\emph{not} vanish at a polynomial rate but is instead controlled by the
regularity of $\tilde{g}$ away from the singularity and the concentration of the
PG$(1,0)$ measure near zero.

The total bias is therefore the sum of the two quadrature errors:
\[
  \operatorname{Bias}(\hat{I}_\text{PCPG-Quad})
  = \underbrace{\frac{1}{2}\,\E_\gamma[\tilde{g}(\gamma) - g(\gamma)]}_{\text{inner (GH), negligible for }N_H \geq 20}
  +\; \underbrace{\frac{1}{2}\!\left(\sum_{j=1}^{M} w_j\,\tilde{g}(\gamma_j)
     - \E_\gamma[\tilde{g}(\gamma)]\right)}_{\text{outer (Gauss--PG), decreasing in }M}.
\]

\paragraph{Bias--variance tradeoff.}
PCPG-Quad has zero variance and a bias that decreases as $M$ and $N_H$ increase.
For a fixed evaluation budget of $M \times N_H$ CF calls, the optimal strategy is
to increase $N_H$ until the inner bias is negligible (cheaply achieved at $N_H \approx 30$)
and then increase $M$ to drive down the outer bias.
The hard limit $M \leq 19$ (imposed by numerical precision in the Golub--Welsch
construction; see \S\ref{subsec:pg-quadrature}) caps the accuracy of
PCPG-Quad regardless of $N_H$.
In practice, PCPG-Quad with $M = 15$--$19$ and $N_H = 30$ gives a highly
accurate, zero-variance estimate that can serve as a reference value against
which the stochastic estimators are validated.

\subsection{Budget Comparison}

To compare all estimators on equal footing we define a shared \emph{total budget}
$N$ as the total number of characteristic-function evaluations:
$N = n_\text{samples}$ for MC,
$N = n_{pg} \cdot n_\text{gauss}$ for PCPG, and
$N = n_{pg} \cdot N_H$ for PCPG-GH.
For PCPG the natural equal-split is $n_{pg} = n_\text{gauss} = \sqrt{N}$;
for PCPG-GH, fixing $N_H$ and setting $n_{pg} = N / N_H$ gives:
\[
  \operatorname{Var}(\hat{I}_\text{PCPG-GH}) \approx \frac{\sigma^2_B \, N_H}{4\,N}.
\]
This recovers the $1/N$ rate of MC, with constant $\sigma^2_B N_H / 4$ in place of
$\sigma^2_\text{MC}$.  The three estimators thus behave as follows at equal budget $N$:

\[
  \operatorname{Var}(\hat{I}_\text{MC}) = \frac{\sigma^2_\text{MC}}{N}, \qquad
  \operatorname{Var}(\hat{I}_\text{PCPG}) \sim \frac{\sigma^2_B}{4\sqrt{N}}, \qquad
  \operatorname{Var}(\hat{I}_\text{PCPG-GH}) \approx \frac{\sigma^2_B N_H}{4N}.
\]

PCPG (full MC) decays as $1/\sqrt{N}$ and is therefore the least efficient of the
three.  PCPG-GH and MC both decay as $1/N$; PCPG-GH is preferable to MC when
$\sigma^2_B N_H < 4\sigma^2_\text{MC}$, i.e.\ when the outer PG variance is small
relative to the sigmoid variance, but in practice $\sigma^2_B$ is typically
comparable to $\sigma^2_\text{MC}$ so MC remains more efficient per evaluation.
The true advantage of the PCPG family over MC is not variance but applicability:
PCPG requires only the characteristic function of $p(\xm \mid \xo)$, not the
ability to draw conditional samples.
A natural equal-split allocation for PCPG sets $n_{pg} = n_\text{gauss} = \sqrt{N}$, giving
\[
  \operatorname{Var}(\hat{I}_\text{PCPG}) = \frac{1}{4}\left[
    \frac{\sigma^2_B}{\sqrt{N}} + \frac{\sigma^2_A}{N}
  \right]
  \;\underset{N \to \infty}{\sim}\; \frac{\sigma^2_B}{4\sqrt{N}}.
\]
The dominant term decays as $1/\sqrt{N}$, whereas the MC variance decays as $1/N$.
PCPG is therefore \emph{less} variance-efficient than MC at equal budget under this allocation.

\subsection{Values Outside \texorpdfstring{$[0,1]$}{[0,1]}}

The MC estimator is always in $(0,1)$ because $\sigma$ maps into that interval.
PCPG, by contrast, computes the real part of a Fourier-type integral approximated
by Monte Carlo, and individual realisations are \emph{not} constrained to $[0,1]$:
the estimator can return negative values or values greater than one when $n_{pg}$
or $n_\text{gauss}$ is small.
This is not a correctness problem --- the estimator is unbiased and converges to
a value in $[0,1]$ --- but it is a practical nuisance when the output is
interpreted as a probability.

\subsection{When is PCPG Preferable?}

The variance comparison above may suggest that PCPG offers no advantage over MC.
Its value lies elsewhere: PCPG never requires samples from $p(\xm \mid \xo)$.
It only needs the \emph{characteristic function} of the conditional distribution,
which is available in closed form for GMMs and general smooth decomposable PCs.
In settings where drawing conditional samples is expensive or intractable ---
for instance with non-Gaussian leaf distributions or densely connected circuits ---
PCPG remains applicable while plain MC does not.
Furthermore, when the inner integral is handled by Gauss--Hermite quadrature
(the PCPG-GH variant), the inner variance $\sigma^2_A$ is eliminated entirely,
reducing the variance to $\sigma^2_B / (4\,n_{pg})$, which matches the $1/N$ rate
of MC (with $N = n_{pg}$) up to the constant $\sigma^2_B / 4$.

\section{Related Work}

Here, we just focus on works that also try to calculate 

\[
  \mathbb{E}_{P(\xm \mid \xo)} \left[\sigma(w^T x)\right]
\]

but use a different approach.

\subsection{Naive Conformant Learning (NaCL)}

\citet{khosravi2019expect} first point out that this expectation is hard
to compute in general: even with a naive Bayes feature distribution, it
is NP-hard. They also note that the popular shortcut of mean imputation
only matches the expectation when the classifier is linear \emph{and} the
features are independent --- for the sigmoid it is always biased.

Their fix is to pick a feature distribution on which the expectation
\emph{is} cheap, namely a naive Bayes model over binary features. The
trick is that naive Bayes and logistic regression are dual to each other:
every naive Bayes induces a logistic regression, and many naive Bayes
models induce the same one. They call a naive Bayes \textit{conformant}
with the given LR if both give the same class probabilities on fully
observed inputs, and learn such a distribution via geometric programming
(``naive conformant learning''). Once this naive Bayes is fitted, the
expected prediction is just a marginal inference query and runs in linear
time.

\medskip
\noindent\textbf{Difference to us:} they keep the feature distribution
simple (naive Bayes on binary inputs) so the expectation is exact and
cheap; we keep the feature distribution rich (a GMM or PC on continuous
inputs) and instead use a stochastic trick to handle the sigmoid. They
also have to refit the feature distribution whenever the classifier
changes, while ours is trained once and reused.

\subsection{Joint Modelling via SAEM}

\citet{jiang2020logreg} assume a joint Gaussian model on the covariates,
$x \sim \mathcal{N}(\mu, \Sigma)$, together with logistic regression on
top. Their main contribution is on the training side: they fit
$(\mu, \Sigma, \beta)$ jointly from training data that may itself contain
missing values, using a stochastic-approximation EM algorithm with
Metropolis--Hastings updates.

For prediction at test time, they simply do plain Monte Carlo: sample
$x^{m,(s)} \sim p(\xm \mid \xo)$ from the fitted Gaussian and average,
\[
  p(y \mid \xo) \approx \frac{1}{S} \sum_{s=1}^{S} p\!\left(y \mid \xo, x^{m,(s)}\right)
\]
This is exactly the MC baseline we compare against, just with a single
Gaussian instead of a mixture as the density.

\medskip
\noindent\textbf{Difference to us:} they sample the missing features and
average the sigmoid; we never sample the missing features, but instead
rewrite the sigmoid so that only the characteristic function of
$p(\xm \mid \xo)$ is needed. Their feature distribution is also limited
to a single Gaussian, while ours can capture multimodal shapes.

\section{Why Monte Carlo Alone Is Not Enough}
\label{sec:why-not-mc}

Monte Carlo estimation of $p(y \mid \xo)$ --- drawing $N$ samples
$x^{m,(i)} \sim p(\xm \mid \xo)$ and averaging the sigmoid predictions --- is
unbiased and simple.  However, it has several limitations that become problematic
in the settings of interest.

\subsection{Conditional Sampling May Be Intractable}

MC requires the ability to \emph{draw samples} from $p(\xm \mid \xo)$.
For a GMM with diagonal covariances this is straightforward, but for a general
smooth decomposable PC it is not always tractable.
Computing the conditional \emph{probability} $p(\xo)$ takes linear time in the
circuit size, but drawing a sample from $p(\xm \mid \xo)$ requires a top-down
ancestral pass that branches on sum-node weights recomputed with the evidence.
For more expressive circuit architectures this pass may be infeasible.

By contrast, the \emph{characteristic function} $\phi_{\xm \mid \xo}(s)$ is
always tractable for smooth decomposable PCs: it is an expectation of a function
that factorises over disjoint variable scopes, which PCs are built to handle
efficiently (see \S\ref{subsec:cf-general-pc}).
PCPG therefore applies wherever MC does not.

\subsection{Variance Grows with Prediction Uncertainty and Missingness}

The variance of the MC estimator is $\sigma^2_\text{MC}/N$, where
\[
  \sigma^2_\text{MC} = \operatorname{Var}_{x \sim p(x \mid \xo)}\!\left[\sigma(w^T x)\right].
\]
This quantity is not under our control: it is determined by the problem geometry.
Two effects make it large in practice.

\paragraph{Proximity to the decision boundary.}
The sigmoid $\sigma(a)$ has variance (over Bernoulli draws) $\sigma(a)(1-\sigma(a))$,
which is maximised at $a = 0$ (the decision boundary) and equals $\tfrac{1}{4}$
there.  Whenever the conditional distribution $p(\xm \mid \xo)$ places
substantial mass on both sides of the hyperplane $w^T x = 0$, the integrand
$\sigma(w^T x)$ fluctuates close to its maximum range $[0,1]$, pushing
$\sigma^2_\text{MC}$ towards $\tfrac{1}{4}$.
To achieve a standard deviation of $\varepsilon$, one then needs
$N \geq 1/(4\varepsilon^2)$ samples --- for $\varepsilon = 0.01$ this is
$N = 2500$, and for $\varepsilon = 0.001$ it is $N = 250{,}000$.

\paragraph{Number of missing features.}
The linear predictor on the missing coordinates is $w_m^T \xm$, a random
variable with variance $w_m^T \Sigma_{\xm \mid \xo} w_m$ under the conditional.
As the number of missing features $d_m = |\{j : x_j \text{ missing}\}|$ grows,
this variance generally increases (more uncertainty is integrated out), which
shifts the distribution of $\sigma(w^T x)$ away from the extremes $\{0, 1\}$
towards the centre, increasing $\sigma^2_\text{MC}$.
In a high-missingness regime --- common in clinical data, sensor dropout, or
survey non-response --- the effective number of samples required for a
fixed precision can be orders of magnitude larger than in the fully-observed
setting.

\subsection{Gradients Through the Estimate Are Noisy}

If one wishes to \emph{learn} the density model $p_\theta(x)$ jointly with
the prediction objective --- for instance, by maximising the log-likelihood of
$p(y \mid \xo)$ over a training set --- then one needs to differentiate
$\E_{p_\theta(\xm \mid \xo)}[\sigma(w^T x)]$ with respect to $\theta$.
For an MC estimator, this requires either:
\begin{itemize}
  \item the \emph{reparameterisation trick}: $x = f(\theta, \varepsilon)$ for
        some noise $\varepsilon$ independent of $\theta$, which is not always
        available for structured or discrete distributions; or
  \item the \emph{score-function (REINFORCE) estimator}:
        $\nabla_\theta \E[f(x)] = \E[f(x)\,\nabla_\theta \log p_\theta(x)]$,
        which is unbiased but notoriously high-variance.
\end{itemize}
The PCPG estimator, by contrast, evaluates the conditional CF
$\phi_{\xm \mid \xo}(s; \theta)$, which is a smooth, differentiable function
of $\theta$ for any parametric density model.  Gradients of the PCPG estimate
with respect to $\theta$ are therefore available by automatic differentiation
without any additional variance-reduction machinery.

\subsection{No Closed Form Exists for the Sigmoid Integral}

Even for the simplest non-trivial density --- a single Gaussian
$p(\xm \mid \xo) = \mathcal{N}(\mu, \Sigma)$ --- the integral
\[
  \E\!\left[\sigma(w^T x) \mid \xo\right]
  = \int_{-\infty}^{\infty} \sigma(\mu_w + \sigma_w z)\,
    \frac{e^{-z^2/2}}{\sqrt{2\pi}}\,\mathrm{d}z,
  \qquad \mu_w = w_m^T\mu,\quad \sigma_w^2 = w_m^T \Sigma w_m,
\]
has no closed form.  Approximations exist (e.g.\ replacing $\sigma$ by the
probit function $\Phi$), but they introduce systematic errors whose magnitude
is problem-dependent.  For a GMM the approximation must be applied per
component and the results combined, compounding the error.  This means that
any accurate, unbiased estimator of the integral must be numerical.

\subsection{Summary}

Table~\ref{tab:method-comparison} summarises how MC and PCPG compare across
these dimensions.

\begin{table}[h]
\centering
\caption{MC vs.\ PCPG for computing $\E[\sigma(w^T x) \mid \xo]$.}
\label{tab:method-comparison}
\begin{tabular}{lcc}
\hline
& \textbf{Monte Carlo} & \textbf{PCPG (ours)} \\
\hline
Unbiased                             & Yes      & Yes (in limit) \\
Requires sampling from $p(\xm \mid \xo)$ & Yes  & No \\
Requires CF of $p(\xm \mid \xo)$     & No       & Yes \\
Variance at budget $N$               & $\sigma^2_\text{MC}/N$ & $\sigma^2_B N_H/(4N)$\rlap{$^*$} \\
Gradients w.r.t.\ density parameters & Noisy    & Automatic differentiation \\
Applicable to general PCs            & Limited  & Yes \\
\hline
\multicolumn{3}{l}{${}^*$PCPG-GH with $n_{pg} = N/N_H$; full-MC PCPG has additional inner variance.}
\end{tabular}
\end{table}

